\section{Số gần đúng và sai số}

\subsection{Đề 1}
\Opensolutionfile{ans}[ans/ans-0-B1]
\caulc
\begin{ex}%[0D6B1-3]
	Khi sử dụng máy tính bỏ túi với $ 10$ chữ số thập phân ta được $\sqrt{8}=2,828427125$. Giá trị gần đúng của $\sqrt{8}$ chính xác đến hàng phần trăm là
	\choice
	{$ 2{,}81$}
	{\True $ 2{,}83$}
	{$ 2{,}82$}
	{$ 2{,}80$}
	\loigiai{
		$2,828427125 \approx 2{,}83$.
	}
\end{ex}

\begin{ex}%%[0D6B1-3]
	Số quy tròn của của $20 \, 182 \, 020$ đến hàng trăm là
	\choice
	{\True $20\, 182 \, 000$}
	{$20 \, 180 \, 000$}
	{$20\, 182 \, 100$}
	{$20\, 182\, 020$}
	\loigiai{$20 \, 182 \, 020$ $\approx 20\, 182 \, 000$. 
	}
\end{ex}

\begin{ex}%%[0D6B1-4]
	Cho số gần đúng $a=8\,141\,378$ với độ chính xác $d=300$. Hãy viết quy tròn số $a$.
	\choice
	{$8\,141\,400$}
	{$8\,142\,400$}
	{\True $8\,141\,000$}
	{$8\,141\,300$}
	\loigiai{ $d=300$ nên $a$ quy tròn đến hàng nghìn. Khi đó	$a=8\,141\,378 \approx 8\,141\,000$.
	}
\end{ex}

\begin{ex}%[0D6B1-4]
	Số quy tròn của số gần đúng $ a=35{,}675$ với độ chính xác $ 0{,}02$ là
	\choice
	{$35{,}6$}
	{$35{,}67$}
	{$36$}
	{\True $35{,}7$}
	\loigiai{ Với độ chính xác $ 0{,}02$  thì $a$ quy tròn đến phần chục. Khi đó	$a=35{,}675  \approx 35{,}7$.
	}
\end{ex}

\begin{ex}%[0D6B1-4]
	Trong một cuộc điều tra dân số, người ta báo dân số của tỉnh A là $ 1\, 279 \, 425\pm 300$ người. Hãy viết số quy tròn số dân trên.
	\choice
	{$ 1\, 270\, 000$ người}
	{$ 1\, 279\, 400$ người}
	{$ 1\, 280\, 000$ người}
	{\True $ 1\, 279\, 000$ người}
	\loigiai{ $d=300$ nên dân số tỉnh A quy tròn đến hàng nghìn. Khi đó
		$1\, 279\, 425\pm 300 \approx 1\, 279\, 000$.}
\end{ex}

\begin{ex}%%[0D6B1-3]
	Số quy tròn đến hàng phần nghìn của số $ a=0{,}1234$ là
	\choice
	{$ 0{,}124$}
	{$ 0{,}12$}
	{\True $ 0{,}123$}
	{$ 0{,}13$}
	\loigiai{
		$ a=0{,}1234 \approx  0{,}123$.}
\end{ex}

\begin{ex}%%[0D6B1-4]
	Cho giá trị gần đúng của $\pi $ là $ a=3,141592653589$ với độ chính xác $10^{-10}$. Hãy viết số quy tròn của số $a$.
	\choice
	{$ a=3{,}1415926535$}
	{$ a=3{,}1415926536$}
	{$ a=3{,}141592653$}
	{\True $ a=3{,}141592654$}
	\loigiai{ Với độ chính xác $10^{-10}$ thì $a$ quy tròn đến vị trí số thứ $10^{-9}$. Khi đó \\
		$ a=3{,}141592653589  \approx a=3{,}141592654$.}
\end{ex}

\begin{ex}%[0D6B1-4]
	Theo thống kê, dân số Việt Nam năm $ 2002$ là $ 79\, 715 \, 675$ người. Giả sử sai số tuyệt đối của số liệu thống kê này nhỏ hơn $ 10\, 000$ người. Hãy viết số quy tròn của số trên.
	\choice
	{$ 79\, 710\, 000$ người}
	{$ 79\, 716\, 000$ người}
	{\True $ 79\, 720\, 000$ người}
	{$ 79\, 000\, 000$ người}
	\loigiai{
		Do sai số tuyệt đối của thống kê này nhỏ hơn $ 10\, 000$ người nên ta quy tròn số trên đến hàng chục nghìn. Vậy số quy tròn của số $ 79 \, 715 \, 675$ là $ 79\, 720\, 000$.}
\end{ex}

\begin{ex}%[0D6B1-4]
	Đo độ cao một ngọn cây là $h=17{,}14\,\text{m}\pm 0{,}3\,\text{m}$. Hãy viết số quy tròn của số $17{,}14$?
	\choice
	{$17{,}1$}
	{$17{,}15$}
	{$17{,}2$}
	{\True $17$}
	\loigiai{ Do độ chính xác là $0{,}3$ nên $h$ quy tròn đến hàng đơn vị, do đó $h=17{,}14 \approx 17$.
	}
\end{ex}

\begin{ex}%[0D6B1-4]
	Cho số $\overline{a}=4{,}1356\pm 0{,}001$. Số quy tròn của số gần đúng $ 4{,}1356$ là
	\choice
	{$ 4{,}135$}
	{$ 4{,}13$}
	{$ 4{,}136$}
	{\True $ 4{,}14$}
	\loigiai{
		Vì độ chính xác đến hàng phần nghìn (độ chính xác là $ 0{,}001$) nên ta quy tròn số $ 4{,}1356$ đến hàng phần phần trăm theo quy tắc làm tròn. Vậy số quy tròn của số $ 4{,}1356$ là $ 4{,}14$.}
\end{ex}

\begin{ex}%%[0D6B1-2]
	Độ cao của một ngọn núi được ghi lại như sau $\overline{h}=1372{,}5\,\text{m}\pm 0{,}2\,\text{m}$. Độ chính xác $ d$ của phép đo trên là
	\choice
	{$ d=0{,}1\,\text{m}$}
	{$ d=1\,\text{m}$}
	{\True $ d=0{,}2\,\text{m}$}
	{$ d=2\,\text{m}$}
	\loigiai{
		Độ chính xác $ d=0{,}2\,\text{m}$.}
\end{ex}
\begin{ex}%%[0D6B1-2]
	Chiều cao của một ngọn đồi là $\overline{h}=347{,}13$$\pm 0,2$ m. Độ chính xác $ d$ của phép đo trên là
	\choice
	{$ d=347{,}13$ m}
	{$ 347{,}33$ m}
	{\True $ d=0{,}2$ m}
	{$ d=346{,}93$ m}
	\loigiai{
		Ta có $ a$ là số gần đúng của $\overline{a}$ với độ chính xác $ d$ qui ước viết gọn là $\overline{a}=a\pm d$. Vậy độ chính xác của phép đo là $ d=0{,}2$ m.}
\end{ex}

\Closesolutionfile{ans}
\indapan{6}{ans/ans-0-B1}
\cauds
\Opensolutionfile{ans}[ans/ans-0-B1-DS]

\begin{ex}%[0D6B1-4]
	Một công ty sử dụng dây chuyền A để đóng gạo vào bao với khối lượng mong muốn $5$ kg. Trên bao bì ghi thông tin khối lượng $5\pm 0,2$ kg. Gọi $\bar{a}$ là khối lượng thực tế của một bao gạo do dây chuyền $A$ đóng gói. Khi đó
	\choiceTF
	{\True Số đúng $\bar{a}$}
	{$5{,}2$ là số gần đúng của $\bar{a}$}
	{Độ chính xác là $ d=5$ kg}
	{\True Giá trị của $\bar{a}$ nằm trong đoạn $[4,8;5,2]$}
	\loigiai{
		\begin{itemchoice}
			\itemch Đúng. Vì Số đúng $\bar{a}$, cũng là khối lượng thực tế của bao gạo (chưa biết chính xác $\bar{a}$). 
			\itemch Sai. Vì Khối lượng mong muốn là $5$ kg nên số $5$ là số gần đúng của $\bar{a}$. 
			\itemch Sai. Vì Độ chính xác là $d=0,2$ kg.
			\itemch Đúng. Vì Giá trị $\bar{a}$ nằm trong đoạn $[5-0,2 ; 5+0,2]=[4,8 ; 5,2]$.
		\end{itemchoice}
	}
\end{ex}


\begin{ex}%[0D6B1-1]
	Kết quả đo chiều dài của một cây cầu được ghi $152$ m $\pm$ $0{,}2$ m; kết quả đo chiều cao của một ngôi nhà được ghi là $15{,}2$ m $\pm 0{,}1$ m. Khi đó	
	\choiceTF
	{\True Sai số tương đối trong cách ghi thứ nhất (đo chiều dài của một cây cầu) là $\delta_1\leqslant\dfrac{d_1}{\left|a_1\right|}=\dfrac{0{,}2}{152}\approx 0,13\%$}
	{\True Sai số tương đối trong cách ghi thứ hai (đo chiều cao của một ngôi nhà) là $\delta_2\leqslant\dfrac{d_2}{\left|a_2\right|}=\dfrac{0{,}1}{15{,}2}\approx 0,66\%$}
	{Sai số tương đối trong cách ghi thứ hai (đo chiều cao của một ngôi nhà) lớn hơn $0{,}66\%$}
	{Cách ghi thứ nhất (đo chiều dài cây cầu) có độ chính xác thấp hơn cách ghi thứ hai (đo chiều cao ngôi nhà)}
	\loigiai{
		\begin{itemchoice}
			\itemch Đúng. Vì Sai số tương đối trong cách ghi thứ nhất là $\delta_1\leqslant\dfrac{d_1}{\left|a_1\right|}=\dfrac{0{,}2}{152}\approx 0{,}13\%$ (hay sai số tương đối không vượt quá $0{,}13\%)$.				
			\itemch Đúng. Vì Sai số tương đối trong cách ghi thứ hai là $\delta_2\leqslant\dfrac{d_2}{\left|a_2\right|}=\dfrac{0{,}1}{15{,}2}\approx 0{,}66\%$ (hay sai số tương đối không vượt quá $0{,}66\%$).
			\itemch Sai. Vì Qua đánh giá sai số tương đối trong hai cách ghi, ta thấy $0{,}13\%<0{,}66\%$ nên cách ghi thứ nhất (đo chiều dài cây cầu) có độ chính xác cao hơn cách ghi thứ hai (đo chiều cao ngôi nhà).
			\itemch Sai. Vì Qua đánh giá sai số tương đối trong hai cách ghi, ta thấy $0{,}13\%<0{,}66\%$ nên cách ghi thứ nhất (đo chiều dài cây cầu) có độ chính xác cao hơn cách ghi thứ hai (đo chiều cao ngôi nhà).
		\end{itemchoice}
	}
\end{ex}


\begin{ex}%%[0D6B1-2]
	Kết quả đo chiều dài của một thửa đất là $75{,}4$ m $\pm$ $0{,}5$ m và đo chiều dài của một cây cầu là $466{,}2$ m $\pm$ $0{,}5$ m. Khi đó
	\choiceTF
	{\True Đối với phép đo thửa đất, sai số tương đối không vượt quá $ 0{,}663\%$}
	{\True Đối với phép đo thửa đất, có sai số tương đối là $\dfrac{d}{|a|}=\dfrac{0{,}5}{75{,}4}=\dfrac{5}{754}$}
	{Đối với phép đo chiều dài cây cầu, có sai số tương đối lớn hơn $\dfrac{5}{4662}\approx 0{,}107\%$}
	{\True Phép đo cây cầu có độ chính xác cao hơn phép đo chiều dài của một thửa đất}
	\loigiai{
		\begin{itemchoice}
			\itemch Đúng. Vì Đối với phép đo thửa đất, tỉ số $\dfrac{d}{|a|}=\dfrac{0{,}5}{75{,}4}=\dfrac{5}{754}$ (tức là sai số tương đối không vượt quá $\dfrac{5}{754}\approx 0{,}663\%$).			
			\itemch Đúng. Đúng. Vì Đối với phép đo thửa đất, tỉ số $\dfrac{d}{|a|}=\dfrac{0{,}5}{75{,}4}=\dfrac{5}{754}$ (tức là sai số tương đối không vượt quá $\dfrac{5}{754}\approx 0{,}663\%$).
			\itemch Sai. 	Vì đối với phép đo chiều dài cây cầu, tỉ số $\dfrac{d}{|a|}=\dfrac{0{,}5}{466{,}2}=\dfrac{5}{4662}$ (nghĩa là sai số tương đối không vượt quá $\dfrac{5}{4662}\approx 0{,}107\%$).
			\itemch Đúng. Vì  ta có $\dfrac{5}{754}>\dfrac{5}{4662}$ nên phép đo cây cầu có độ chính xác cao hơn.
		\end{itemchoice}
	}
\end{ex}
\begin{ex}%%[0D6B1-4]
	Xét tính đúng, sai của các mệnh đề sau
	\choiceTF
	{\True Cho số gần đúng $a=2022$ với độ chính xác $d=20$. Số quy tròn của $a$ là $2020$}
	{\True  Cho số gần đúng $a=1{,}04527$ với độ chính xác $d=0{,}004$. Số quy tròn của $a$ là $1{,}05$}
	{Cho $\bar{a}=14{,}6543\pm 0{,}03$. Số quy tròn của $a$ là $14{,}65$}
	{Cho số gần đúng $a=301335$ với độ chính xác $d=234$. Số quy tròn của $a$ là $301000$}
	\loigiai{
		\begin{itemchoice}
			\itemch Đúng. Vì hàng lớn nhất của $d$ là hàng chục nên ta quy tròn $a$ đến hàng trăm. Vậy số quy tròn của $a$ là $2000$. 
			\itemch Đúng. Vì hàng lớn nhất của $d$ là hàng phần nghìn nên ta quy tròn $a$ đến hàng phần trăm. Vậy số quy tròn của $a$ là $1{,}05$.
			\itemch Sai. Vì hàng lớn nhất của $d$ là hàng phần trăm nên ta quy tròn $a$ đến hàng phần mười. Vậy số quy tròn của $a$ là $14{,}7$.
			\itemch Sai. Vì cho số gần đúng $a=301335$ với độ chính xác $d=234$. Kết quả là $301000$.
		\end{itemchoice}
	}
\end{ex}
\Closesolutionfile{ans}
\indapan{3}{ans/ans-0-B1-DS}

\Opensolutionfile{ans}[ans/ans-0-B1-KQ]
\caukq
\begin{ex}%%[0D6B1-2]
	Bạn Nam thực hiện việc đo đường kính của một ống đồng bằng thước kẹp. Kết quả của Nam là $15{,}3$ mm. Hỏi số  số gần đúng của phép đo là bao nhiêu?
	\shortans{$15{,}3$}
	\loigiai{
		Số $15{,}3$ là số gần đúng của phép đo trong thực tế}
\end{ex}
\begin{ex}%%[0D6K1-2]
	Trên bao bì của một sản phẩm có ghi \lq\lq khối lượng tịnh $200\pm 2$ g\rq\rq. Hãy đánh giá sai số tương đối của số gần đúng này.
	\shortans{ $ 1$}
	\loigiai{
		Khối lượng đúng của bao sản phẩm $\bar{a}$ (tính theo gam) thuộc đoạn $[198; 202]$. \\
		Sai số tương đối là  $\delta_a\leq\dfrac{2}{200}=1\%$.}
\end{ex}
\begin{ex}%[0D6B1-1]
	Làm tròn số $2315{,}56$ đến hàng đơn vị. Tính sai số tuyệt đối của chúng.
	\shortans{$0{,}44$}
	\loigiai{
		Số $2315{,}56$ được làm tròn đến hàng đơn vị là $2316$ và sai số tuyệt đối là \\ $\Delta=|2315{,}56-2316|=0{,}44$.
	}
\end{ex}
\begin{ex}%[0D6K1-1]
	Bạn Ngân có một mảnh nhựa với bề mặt hình tròn bán kính $1$ dm. Bạn ấy thực hiện đo chu vi của mép mảnh nhựa đó bằng cách sử dụng một sợi dây dài không dãn như sau: Cố định một đầu sợi dây trên mép mảnh nhựa, rồi quấn sợi dây quanh mép mảnh nhựa một vòng cho đến khi đầu dây cố định chạm vào thân sợi dây lần đầu tiên, sau đó đo độ dài phần dây chạm vào mép mảnh nhựa và được kết quả là $6$ dm. Khi đó sai số tuyệt đối trong phép đo không vượt quá bao nhiêu dm.
	\shortans{ $0{,}3$}
	\loigiai{
		Chu vi của mép mảnh nhựa là: $\bar{a}=2\pi \cdot 1=2\pi$ (dm).\\
		Gọi chu vi của mép mảnh nhựa mà bạn Ngân đo được là $a$, suy ra $a=6$ (dm). Vì $3<\pi<3{,}15$ nên $6<2\pi<6{,}3\Rightarrow 6<\bar{a}<6{,}3\Rightarrow\Delta_a=|\bar{a}-6|<6{,}3-6=0,3$. Suy ra sai số tuyệt đối trong phép đo không vượt quá $0{,}3$ dm.}
\end{ex}

\begin{ex}%[0D6K1-1]
	Bạn Ngân có một mảnh nhựa với bề mặt hình tròn bán kính $1$ dm. Bạn ấy thực hiện đo chu vi của mép mảnh nhựa đó bằng cách sử dụng một sợi dây dài không dãn như sau: Cố định một đầu sợi dây trên mép mảnh nhựa, rồi quấn sợi dây quanh mép mảnh nhựa một vòng cho đến khi đầu dây cố định chạm vào thân sợi dây lần đầu tiên, sau đó đo độ dài phần dây chạm vào mép mảnh nhựa và được kết quả là $6$ dm. Khi đó sai số tương đối trong phép đo không vượt quá bao nhiêu $\%$.
	\shortans{ $5$}
	\loigiai{
		Ta có $\Delta_a<0{,}3$ nên $\dfrac{\Delta_a}{|a|}<\dfrac{0{,}3}{6}=0{,}05=5\%$. Suy ra sai số tương đối trong phép đo không vượt quá $5\%$.}
\end{ex}
\begin{ex}%Câu 16
	Dùng phân số $\dfrac{33}{19}$ để làm số gần đúng cho $\sqrt{3}$. Hãy đánh giá sai số tuyệt đối mắc phải là bao nhiêu?
	\shortans{$ 0{,}01$}
	\loigiai{
		Sử dụng máy tính cầm tay, ta có  $\sqrt{3}=1{,}732050808\ldots ;\dfrac{33}{19}=1{,}736842105\ldots$\\
		Sai số tuyệt đối của số gần đúng là\\
		$\begin{aligned}
			\left|\sqrt{3}-\dfrac{33}{19}\right|& =|1,732050808\ldots-1{,}736842105\ldots |\\ 
			&=1{,}736842105\ldots-1{,}732050808\ldots <1{,}74-1{,}73=0{,}01.\\ 
		\end{aligned}$\\
	}
\end{ex}
\Closesolutionfile{ans}
\indapan{6}{ans/ans-0-B1-KQ}
\newpage 
\subsection{Đề 2}
\Opensolutionfile{ans}[ans/ans-0-B1]
\caulc

%17
\begin{ex}%[0D6H1-3]
	Cho giá trị gần đúng của $\pi$ là $a=3,141\,592\,653\,589$ với độ chính xác $10^{-10}$ ($10$ chữ số thập phân). Số quy tròn của $a$ là
	\choice
	{\True $a=3{,}141592654$}
	{$a=3{,}1415926536$}
	{$a=3{,}141592653$}
	{$a=3{,}1415926535$}
	\loigiai{
		Ta có $10^{-11}<10^{-10}<10^{-9}$ nên hàng cao nhất mà $d$ nhỏ hơn một đơn vị của hàng đó là hàng phần tỉ. Do đó ta phải quy tròn số $a=3{,}141\,592\,653\,589$ đến hàng phần tỉ.\\
		Vậy, số quy tròn là $a= 3{,}141\,592\,654$.}
\end{ex}
%câu 26
\begin{ex}%[0D6H1-3]
	Cho số $\overline{a}=17\,658\pm 16$. Số quy tròn của số gần đúng $17\,658$ là
	\choice
	{$18\,000$}
	{$17\,800$}
	{$17\,600$}
	{\True $17\,700$}
	\loigiai{
		Ta có $10<16<100$ nên hàng cao nhất mà $d$ nhỏ hơn một đơn vị của hàng đó là hàng trăm. Do đó ta phải quy tròn số $176\,38$ đến hàng trăm.\\
		Vậy số quy tròn là $17\,700$ (hay viết $\overline{a}\approx 17\,700$).
	}
\end{ex}
%19
\begin{ex}%[0D6H1-3]
	Cho $\overline{a}=31\,462\,689\pm 150$. Số quy tròn của số $31\,462\,689$ là
	\choice
	{$31\,462\,000$}
	{$31\,463\,700$}
	{$31\,463\,600$}
	{\True $31\,463\,000$}
	\loigiai{
		Độ chính xác đến hàng trăm $(d=150)$ nên ta quy tròn đến hàng nghìn. Vậy số quy tròn của số $31\,462\,689$ là $31\,463\,000$.}
\end{ex}
%20
\begin{ex}%[0D6H1-3]
	Theo thống kê, dân số Việt Nam năm $2016$ được ghi lại như sau $\overline{S}=94\,444\,200\pm 3000$ (người). Số quy tròn của số gần đúng $94\,444\,200$ là
	\choice
	{\True $94\,440\,000$}
	{$94\,450\,000$}
	{$94\,444\,000$}
	{$94\,400\,000$}
	\loigiai{
		Vì $1000<3000<10000$ nên hàng cao nhất mà $d$ nhỏ hơn một đơn vị của hàng đó là hàng chục nghìn. Nên ta phải quy tròn số $94\,444\,200$ đến hàng chục nghìn. Vậy số quy tròn là $94\,440\,000$.}
\end{ex}
%21
\begin{ex}%[0D6H1-3]
	Cho số $a=367\,653\,964\pm 213$. Số quy tròn của số gần đúng $367\,653\,964$ là
	\choice
	{$367\,653\,960$}
	{$36\,7653\,000$}
	{\True $367\,654\,000$}
	{$367\,653\,970$}
	\loigiai{
		Vì độ chính xác đến hàng trăm $d=213$ nên số quy tròn của số gần đúng $367\,653\,964$ là $36\,7654\,000$.}
\end{ex}
%22
\begin{ex}%[0D6H1-3]
	Theo thống kê, dân số Việt Nam năm $2002$ là $79\,715\,675$ người. Giả sử sai số tuyệt đối của số liệu thống kê này nhỏ hơn $10\,000$ người. Hãy viết số quy tròn của số trên
	\choice
	{$79\,710\,000$ người}
	{$79\,716\,000$ người}
	{\True $79\,720\,000$ người}
	{$79\,700\,000$ người}
	\loigiai{
		Vì sai số tuyệt đối của số liệu thống kê này nhỏ hơn $10\,000$ người nên độ chính xác đến hàng nghìn nên ta quy tròn đến hàng chục nghìn.\\
		Vậy số quy tròn cùa số trên là $79\,720\,000$ người.
	}
\end{ex}
%23
\begin{ex}%[0D6H1-3]
	Độ dài các cạnh của một đám vườn hình chữ nhật là $x=7{,}8$ $\pm 2$ cm và $y=25{,}6$ m $\pm 4$ cm. Cách viết chuẩn của diện tích (sau khi quy tròn) là
	\choice
	{$200\mathrm{~m}^2\pm 0{,}9\mathrm{~m}^2$}
	{$199\mathrm{~m}^2\pm 0{,}8\mathrm{~m}^2$}
	{$199\mathrm{~m}^2\pm 1\mathrm{~m}^2$}
	{\True $200\mathrm{~m}^2\pm 1\mathrm{~m}^2$}
	\loigiai{
		Ta có 
		\begin{eqnarray*}
			x=7{,}8 \text{ m} \pm 2 \text{ cm} &\Rightarrow& 7{,}78 \text{ m}\leq x\leq 7{,}82 \text{ m}.\\
			y=25{,}6 \text{ m} \pm 4 \text{ cm} &\Rightarrow& 25{,}56 \text{ m} \leq y\leq 25{,}64 \text{ m}.
		\end{eqnarray*}
		Do đó diện tích của hình chữ nhật thỏa mãn $198{,}8568\mathrm{~m}^2\leq xy\leq 200{,}5048\mathrm{~m}^2$.\\
		Vậy, cách viết chuẩn của diện tích sau khi quy tròn là $200$ m$^2\pm 1$ m$^2$.
	}
\end{ex}
%câu 24
\begin{ex}%[0D6H1-3]
	Chiều dài gần đúng của một cái bàn học là $a=1{,}238$ m với độ chính xác $d=0{,}01$ m. Hãy viết số quy tròn của số $a$?
	\choice
	{\True Số quy tròn của $a$ là $1{,}2$ m}
	{Số quy tròn của $a$ là $1{,}23$ m}
	{Số quy tròn của $a$ là $1{,}24$ m}
	{Số quy tròn của $a$ là $1{,}248$ m}
	\loigiai{
		Vì độ chính xác $d=0{,}01$ đến hàng phần trăm nên ta quy tròn số gần đúng $a$ đến hàng phần chục.\\
		Vậy, số quy tròn của $a$ là $1{,}2$ m.
	}
\end{ex}
%câu 25
\begin{ex}%[0D6H1-3]
	Số tiền quỹ lớp $10$A còn lại là $a=1\,647\,500$ (đồng) với độ chính xác $d=500$ (đồng). Hãy viết số quy tròn của số $a$?
	\choice
	{\True Số quy tròn của $a$ là $1\,648\,000$ (đồng)}
	{Số quy tròn của $a$ là $1\,647\,000$ (đồng)}
	{Số quy tròn của $a$ là $1\,649\,000$ (đồng)}
	{Số quy tròn của $a$ là $1\,650\,000$ (đồng)}
	\loigiai{
		Vì độ chính xác  $d=500$ đến hàng trăm nên ta quy tròn số gần đúng $a$ đến hàng nghìn.\\
		Vậy, số quy tròn của $a$ là $1\,648\,000$  (đồng).	
	}
\end{ex}

%27
\begin{ex}%[0D6H1-3]
	Cho số $\overline{a}=4{,}1356 \pm 0{,}001$. Số quy tròn của số gần đúng $4{,}1356$ là
	\choice
	{ $4{,}135$}
	{$4{,}13$}
	{$4{,}136$}
	{\True $4{,}14$}
	\loigiai{
		Vì độ chính xác đến hàng phần nghìn (độ chính xác là $0{,}001$) nên ta quy tròn số $4{,}1356$ đến hàng phần phần trăm theo quy tắc làm tròn. Vậy số quy tròn của số $4{,}1356$ là $4{,}14$.}
\end{ex}	
%28
\begin{ex}%%[0D6H1-1]
	Cho giá trị gần đúng của $\dfrac{8}{17}$ là $0{,}47$. Sai số tuyệt đối của $0{,}47$ là
	\choice
	{\True $0{,}001$}
	{$0{,}002$}
	{$0{,}003$}
	{$0{,}004$}
	\loigiai{
		Ta có $\dfrac{8}{17}=0{,}470588235294 \ldots$.\\
		Sai số tuyệt đối của $0{,}47$ là $\left|0{,}47-\dfrac{8}{17}\right|<|0{,}47-0{,}471|=0{,}001$.
	}
\end{ex}	
%29
\begin{ex}%[0D6H1-1]
	Đo chiều dài của một cây thước, ta được kết quả $\overline{a}=45 \pm 0{,}3$ cm. Khi đó sai số tuyệt đối của phép đo được ước lượng là
	\choice
	{$\Delta_{45}=0{,}3$}
	{\True $\Delta_{45} \leq 0{,}3$}
	{$\Delta_{45} \leq -0{,}3$}
	{$\Delta_{45} = 0{,}3$}
	\loigiai{
		Ta có độ dài dài gần đúng của cây thước là $a=45$ với độ chính xác $d=0{,}3$.\\	
		Nên sai số tuyệt đối $\Delta_{45} \leq d=0{,}3$
		.}
\end{ex}	

\Closesolutionfile{ans}
\indapan{6}{ans/ans-0-B1}
\cauds
\Opensolutionfile{ans}[ans/ans-0-B1-DS]

%câu 6
\begin{ex}%[0D6V1-4]
	Các mệnh đề sau đúng hay sai?
	\choiceTF
	{Cho số gần đúng $a=581\,268$ với độ chính xác $d=200$. Có số quy tròn là $581\,200$}
	{\True Cho số gần đúng của $\pi$ là $a=3{,}141592653589$, độ chính xác là $10^{-10}$. Số quy tròn của $a$ là $3{,}141592654$}
	{\True Chiều dài một cái cầu đo được là $l=1745{,}25$ $\pm 0{,}01$ m. Có số quy tròn là $1745{,}3$ m}
	{\True Số gần đúng $\sqrt{5}$ với độ chính xác $0{,}005$ là $\approx 2{,}24$}
	\loigiai{
		\begin{itemchoice}
			\itemch Sai. Vì độ chính xác được cho đến hàng trăm $(d=200)$ nên ta cần quy tròn số gần đúng đến hàng nghìn. Do đó ta thu được số quy tròn là $581\,000$. 
			\itemch Đúng. Vì độ chính xác của số gần đúng đến $10^{-10}$ ($10$ chữ số thập phân sau dấu phẩy) nên ta quy tròn số đó đến $10^{-9}$ ($9$ chữ số thập phân sau dấu phẩy).
			Vậy số quy tròn của $a$ là $3{,}141 592654$.
			\itemch Đúng. Ta có  $l=1745{,}25$ m $\pm 0,01$ m có độ chính xác đến hàng phần trăm (độ chính xác là $0{,}01$) nên ta quy tròn số gần đúng đến hàng phần chục.
			Vậy số quy tròn của $1745{,}25$ m đến hàng phần chục là $1745{,}3$ m.
			\itemch Đúng. Số gần đúng $\sqrt{5}$ với độ chính xác $0{,}005$ là $\approx 2{,}24$.
		\end{itemchoice}
	}
\end{ex}
%câu 7
\begin{ex}%[0D6H1-3]
	Xét tính đúng, sai của các mệnh đề sau
	\choiceTF
	{\True $230\,592\pm 300$ có số quy tròn là $231\,000$}
	{$1{,}62516248\pm 0{,}001$ có số quy tròn là $1{,}635$}
	{\True $6\,491\,258\pm 1000$ có số quy tròn là $6\,490\,000$}
	{\True $564{,}65449\pm 0{,}003$ có số quy tròn là $564{,}65$}
	\loigiai{
		\begin{itemchoice}
			\itemch Đúng. $231\,000$. 
			\itemch Sai. $1{,}63$.
			\itemch Đúng. $6\,490\,000$.
			\itemch Đúng. $564{,}65$.
		\end{itemchoice}
	}
\end{ex}
%câu 8
\begin{ex}%[0D6V1-4]
	Xét tính đúng, sai của các mệnh đề sau
	\choiceTF
	{Cho $\overline{a}=\sqrt{7}=2{,}645751$. Số gần đúng của $\overline{a}$ với độ chính xác $d=0{,}004$ là $2{,}6$}
	{Cho $\overline{a}=\dfrac{\sqrt[3]{2}}{4}=0{,}314980\ldots$ Số gần đúng của $\overline{a}$ với độ chính xác $d=0{,}01$ là $0{,}315$}
	{\True Cho số gần đúng $a=1\,624\,192$ với độ chính xác $d=300$. Số quy tròn của $a$ là $1\,624\,000$}
	{\True Viết quy tròn số gần đúng $a=3{,}26356$ biết $\overline{a}=3{,}26356\pm 0{,}001$ là $3{,}26$}
	\loigiai{
		\begin{itemchoice}
			\itemch Sai. Kết quả là $2{,}65$. 
			\itemch Sai. Kết quả là $0{,}3$.
			\itemch Đúng. Số quy tròn của $a$ là $1\,624\,000$.
			\itemch Đúng. $\overline{a}=3{,}26356\pm 0{,}001\Rightarrow$ Độ chính xác đến hàng phần nghìn (độ chính xác là $d=0{,}001$).\\
			Ta quy tròn đến hàng phần trăm. Vậy số quy tròn của $a$ là $3{,}26$.
		\end{itemchoice}
	}
\end{ex}
%câu 9 
\begin{ex}%[0D6V1-4]
	Cho ba giá trị gần đúng của $\dfrac{3}{7}$ là $0{,}429$; $0{,}4$ và $0{,}42$. Xét tính đúng, sai của các mệnh đề sau
	\choiceTF
	{\True Công thức đánh giá sai số tuyệt đối là $\Delta=|\overline{a} - a|$}
	{\True Xét số gần đúng $0{,}429$ ta có $\Delta_1=\left|\dfrac{3}{7} - 0{,}429\right|<0{,}0005$}
	{\True Xét số gần đúng $0{,}4$ ta có $\Delta_2=\left|\dfrac{3}{7} - 0{,}4\right|<0{,}03$}
	{\True Xét số gần đúng $0{,}42$ ta có $\Delta_2=\left|\dfrac{3}{7} - 0{,}42\right|<0{,}009$}
	\loigiai{
		\begin{itemchoice}
			\itemch Đúng. Ta sử dụng công thức đánh giá sai số tuyệt đối là $\Delta=|\overline{a} - a|$.
			\itemch Đúng. Xét số gần đúng $0{,}429$ ta có $\Delta_1=\left|\dfrac{3}{7} - 0{,}429\right|<0{,}0005$.
			\itemch Đúng. Xét số gần đúng $0{,}4$ ta có $\Delta_2=\left|\dfrac{3}{7} - 0{,}4\right|<0{,}03$.
			\itemch Đúng. Xét số gần đúng $0{,}42$ ta có $\Delta_2=\left|\dfrac{3}{7} - 0{,}42\right|<0{,}009$.
		\end{itemchoice}
	}
\end{ex}
%câu 10
\begin{ex}%[0D6V1-4]
	Xét tính đúng, sai của các mệnh đề sau
	\choiceTF
	{\True Số quy tròn của $31416{,}1$ đến hàng chục là $31\,420$}
	{\True Số quy tròn của $31{,}135$ đến hàng phần trăm là $31{,}14$}
	{Số quy tròn của $110{,}32344$ đến hàng phần nghìn là $110{,}323$}
	{\True Giá trị gần đúng của số $\sqrt[3]{2}$ chính xác đến hàng phần trăm là $1{,}26$}
	\loigiai{
		\begin{itemchoice}
			\itemch Đúng. Số quy tròn của $31416{,}1$ đến hàng chục là $31\,420$. 
			\itemch Đúng. Số quy tròn của $31{,}135$ đến hàng phần trăm là $31{,}14$.
			\itemch Đúng. Số quy tròn của $110{,}32344$ đến hàng phần nghìn là $110{,}323$
			\itemch Đúng. Quy tròn số đã cho chính xác đến hàng phần trăm là $1{,}26$.
		\end{itemchoice}
	}
\end{ex}


\Closesolutionfile{ans}
\indapan{3}{ans/ans-0-B1-DS}

\Opensolutionfile{ans}[ans/ans-0-B1-KQ]
\caukq

\begin{ex}%[0D6H1-1]
	Trong một cuộc điều tra dân số, người ta viết dân số của một tỉnh là $3\,574\,625 \pm 50\,000$ (người).
	Hãy đánh giá sai số tương đối của số gần đúng này?
	\shortans{$1{,}4$}
	\loigiai{
		Ta có $a=3\,574\,625$ và $d=50\,000$ nên sai số tương đối $\delta_{a} \leqslant \dfrac{d}{|a|}=\dfrac{50\,000}{3\,574\,625} \approx 1{,}4 \%$.
	}
\end{ex}
%15
\begin{ex}%[0D6H1-1]
	Bạn Lan tính diện tích hình tròn bán kính $r=3$ cm bằng công thức $S=3{,}14 \cdot 3^{2}=28{,}26$ cm$^2$.
	Biết rằng $3{,}1<\pi<3{,}2$, hãy ước lượng sai số tương đối của $S$.
	\shortans{$1{,}91$}
	\loigiai{
		Diện tích đúng kí hiệu là $\overline{S}$ thỏa mãn $3{,}1\cdot 3^{2}<\overline{S}<3{,}2\cdot 3^{2} \Leftrightarrow 27{,}9<\overline{S}<28{,}8$.\\
		Do đó $27{,}9-28{,}26<\overline{S}-S<28{,}8-28{,}26 \Leftrightarrow-0{,}36<\overline{S}-S<0{,}54$.\\
		Suy ra $|\overline{S}-S|<0{,}54$.
		Vậy sai số tương đối không vượt quá $\dfrac{0{,}54}{28{,}26} \approx 1{,}91 \%$.
		
	}
\end{ex}

%17
\begin{ex}%[0D6V1-1]
	Một giá trị gần đúng của $\pi$ là $3{,}142$. Hãy ước lượng sai số tuyệt đối và sai số tương đối của giá trị gần đúng trên biết rằng $3{,}141<\pi<3{,}143$.
	\shortans{$0{,}03$}
	\loigiai{
		Một giá trị gần đúng của $\pi$ là $3{,}142$. Hãy ước lượng sai số tuyệt đối và sai số tương đối của giá trị gần đúng trên biết rằng $3{,}141<\pi<3{,}143$.\\
		Ta có $3{,}141 - 3{,}142<\pi - 3{,}142<3{,}143 - 3{,}142\Leftrightarrow - 0{,}001<\pi - 3{,}142<0{,}001$.\\ Suy ra$\pi - 3{,}142\mid<0{,}001$.\\
		Do đó sai số tuyệt đối là $0{,}001$.
		Sai số tương đối không vượt quá $\dfrac{0,001}{|3,142|}<\dfrac{0,001}{3,142}\approx 0,0318\%$.
	}
\end{ex}
%18
\begin{ex}%[0D6V1-3]
	Hãy quy tròn số $\overline{a}=\dfrac{5}{7}=0{,}714285$ đến hàng phần trăm và ước lượng sai số tương đối.
	\shortans{$0{,}7$}
	\loigiai{
		Hãy quy tròn số $\overline{a}=\dfrac{5}{7}=0{,}714285$ đến hàng phần trăm và ước lượng sai số tương đối.\\
		Kết quả quy tròn số $\overline{a}=\dfrac{5}{7}=0{,}714285$ đến hàng phần trăm là $0{,}71$.\\
		Sai số tương đối không vượt quá $\dfrac{0{,}005}{071}\approx 0{,}70\%$.
	}
\end{ex}

%20
\begin{ex}%[0D6V1-3]
	Cho số gần đúng $a=2\,362$ với độ chính xác $d=100$. Hãy viết số quy tròn của số $a$ và ước lượng sai số tương đối của số quy tròn đó.
	\shortans{$23{,}1$}
	\loigiai{
		Số quy tròn của $a$ là $2\,000$.\\
		Ta có $a-d \leq \overline{a} \leq a+d \Leftrightarrow 2262 \leq \overline{a} \leq 2462$.\\
		$262 \leq \overline{a}-2000 \leq 462 \Leftrightarrow \Delta_{2000}=|\overline{a}-2000| \leq 462$.\\
		Sai số tương đối của $2000$ là $\delta_{2000}=\dfrac{\Delta_{2000}}{|2000|} \leq \dfrac{462}{2000}=23{,}1 \%$.
	}
\end{ex}


%24
\begin{ex}%[0D6V1-3]
	Độ dài các cạnh của mảnh vườn hình chữ nhật là $x=7{,}8$ m $\pm 2$ cm và $y=25{,}6$ m $\pm 4$ cm. Tìm diện tích (sau khi quy tròn) của mảnh vườn.
	\shortans{$200$}
	\loigiai{
		Ta có $\heva{&7{,}78\leq x\leq 7{,}82\\&25{,}56\leq y\leq 25{,}64}\Rightarrow 198{,}8568\leq \overline{S}=xy\leq 200{,}5048$.\\
		Diện tích gần đúng là $S=7{,}8\cdot 25{,}6=199{,}68$.\\
		Suy ra $-0{,}8232\leq \overline{S} - S\leq 0{,}8248\Rightarrow\Delta_S=|\overline{S} - S|\leq 0{,}8248$.
		Suy ra diện tích gần đúng là $S=199{,}68$ m$^2$ với độ chính xác $d=0{,}8248$.
		Nên số quy tròn của $S$ là $200$ m$^2$.
	}
\end{ex}

\Closesolutionfile{ans}
\indapan{6}{ans/ans-0-B1-KQ}